\chapter{影响西藏高中生数学运算能力的因素}


\section{从学生角度分析}

\subsection{运算习惯有待培养}

大部分学生做题前不认真审题，数据抄写错误或看错行，这就是学生经常会出现马虎的重要原因之一；做题时思路不清晰，运算程序混乱；做题后不认真检查；还有很多学生不能规范书写运算过程。这些不良习惯，大大的增加了出错率。


\subsection{对数学基础知识的理解与掌握不到位}

大多数同学对于不同的数学方程没有很好的记住，也没有很好地使用这些方程。记住了不同的方程，却无法正确地选择正确的答案。比如在求解一元二次方程式时，不能记住分母是2；在解决三维立体几何问题时，不能正确地构建或不能在坐标系统中精确定位等问题。事实上，每一次的数学操作都是建立在一些基本的基本方程之上。这些人的基本功都很差，做题的时候没有思考，也没有考虑过这些问题，所以他们的成绩一直很差。


\subsection{对数学不感兴趣
}

在实习期间，我发现许多西藏学生对数学不感兴趣，跟学生平时交流时侯他们一直会说数学好难，他们的数学基础知识不牢固，久而久之导致他们对数学不感兴趣。小学与初中的数学还比较简单，学生还比较有兴趣。高中数学比小学和初中的难度上升了一个大阶梯，数学运算量大且难，且高中科目众多，学生压力大，数学一道难题就需要很长时间解答还不一定能解出来甚至解出来还是错的，耗费了绝大多数同学很大一部分精力，以至于他们没有时间来做其他科目。这也是学生对学习数学兴趣降低的原因之一，也是学生数学运算能力下降的原因之一。



\subsection{对错题没有及时整理、总结、归纳}

反思，总结，归纳是学习数学的重要一环。整理、归纳、归纳是对学生认知的一次发现、总结、归纳，对于加强西藏中学学生的学习、总结、提升学生的数学素养都是非常有用的，整理、整理是提升学生的数学能力的有效途径，然而，出现问题的原因不仅仅是技术上的问题，还必须对其进行整理和整理。然而，因为高中有大量的家庭作业，因此很少有时间去整理和反思那些不正确的问题。然而，教师在教学中却常常注重于教学的内容，忽视了对学生进行思维和总结的学习。课后反思、总结、归纳会导致西藏中学学生形成思维定势，缺乏自我反思和总结。


\subsection{对数学运算认识不到位}

当前，学生对数学运算的认识不足，停留在计算量大、难度大的领域，对数学感到厌烦。课上遇到例题学生只看思路，不会自己重新运算；课后错题也只看思路，对错题从来都归结到马虎。教师忽略了对学生数学运算的认识，导致了学生对数学产生了恐惧，教师教学没有效率，学生学习也起不到好的效果。



\section{从教师角度分析}

\subsection{缺乏对数学运算的培养}

在教学中，老师很少针对学生的数学计算能力进行专门的训练和系统的训练。
教师只有在课堂上才能系统全面的讲授知识，但一节课的时间有限，数学教师在课上只会讲授相关知识或只会在黑板上写下重点公式和关键得分点，而往往省略运算步骤，导致教育工作水平降低。


\subsection{认识不到自己的责任}

许多中学数学老师都觉得，数学计算应该由初中生或者小学生来完成，而老师却偏偏把注意力放在了数学思考上，一头栽到了茫茫的题目里，结果却什么也没有，让老师和同学感到非常的头疼。

比如：在中学里，二次方程、一元一次方程、二次方程、一元一次方程的计算，而在中学里，交叉相加很少见，大部分时间里都是用方程式和计算方程式，所以，他们并没有太好的学习经验\cite{李雅男西藏高中生数学运算能力的培养}。

例如，一元二次方程式$18 x^{2}-21 x+5=0$，用十字相乘法时得到 $(3 {x}-1)(6 {x}-5)=0$, 就可以很容易地得出$X_{1}=\frac{1}{3}, x_{2}=\frac{5}{6}$，但如果使用 $b^{2}-4 {ac} \geqslant 0$ ，则会出现错误。所以，中学数学老师在传授二次功能的时候，必须首先对交叉相乘进行温习，使其应用更加娴熟，从而使其在一定的学习过程中完成更多的练习。


\subsection{教师的旧课堂}

以前的数学课老师总是在讲授，或是直接背诵，让他们完成很多的练习。高中生的课程内容很多，一般的高中生一节课四十五分钟，一天两节，老师们会在规定的时间内教授他们的课程，而老师们则会给学生们更多的自由讨论和互动，从而降低他们的学习效率。很多老师都会因为学生的听课效率和突然的变化而受到干扰，大多数的老师都觉得西藏的学生很有自主的学习和自主，所以上课时忽略了学生的讨论和老师的提问。



\section{从教学环境角度分析}

\subsection{为考试产生的教学}

在高考的时候，教师和家长只注重分数，只注重分数，只注重分数，只注重分数，所以，很多时候，数学题太多，而在课堂上，老师只会给他们讲一些简单的概念，让他们在课堂上完成。另外，在我的实习中，我注意到了我的导师们都是日复一日地学习，所以大部分的训练都是由他们自己来完成。长期下去，会导致学习技能不熟练，操作水平不会得到提升，数学计算能力的培养也会变得困难。


\subsection{新旧知识的衔接}

初中的数学课程不多，学习简单，基本不需要太多的基础，有些知识点课堂上无法理解，但可以在课堂上自行学习。但中学的知识量较大，需要掌握一些基本的知识。中学的某些基础知识和计算技能，在中学里是无法学习和使用的，但是中学老师们觉得，中学老师已经把所有的基础知识都传授给学生了，老师们只关注着课堂上的内容，这样的情况造成了中学与中学的知识之间的脱节。所以，中学教育应该重视初中与中学的结合，这样既能夯实学生的基本知识，又能使学生在中学阶段有一个过渡期。


\subsection{各种辅助工具的出现}

网络、书店和各种学习机在现代社会中的应用对学习既有好处也有坏处。现代社会，电脑等辅助学习的辅助手段的出现，让学习的压力大大减少，但是大部分的同学都过分地依靠辅助学习，觉得这样做既方便快捷又准确，反而会让他们失去学习的积极性和主动性，从而影响运算的效率。市场上各类补习材料良莠不齐，而学生却是瞎买，不知好歹，容易产生误解。网络上的各种教学录像，可以让同学们在课堂上看得清清楚楚，但也有可能导致学生在课堂上不专心，只能靠自学。


\subsection{社会娱乐的吸引}

在这个娱乐方式多种多样的社会，有太多好吃的、好玩的吸引学生的注意。学生在学校完全接触不到这些，所以在每次放假的时候学生都会被这些活动吸引，从而无法专注学习,而开学时又对它们恋恋不舍，很难收心。
